\chapter{Calibration and Performance of SVJJ}
\label{ch:calibration}
To assess the performance of the SVJJ model, we consider a dataset of European call and put options
on the S\&P 500 index (SPX), observed at the close of business on 13 November 2019. On this date,
the SPX closed at 3094.04 USD. The data processing follows \parencite{FEdatahandling} and
\parencite{Kokholm2016}. After removing options with a zero bid or ask price, the dataset contains
3,665 observations across 12 different maturities.

Before filtering the data, we infer the dividend yield $q$ and the continuously compounded spot
rates $r_{T_i}$, $i=1,\ldots,12$. Following \parencite{FEdatahandling}, these are estimated by
solving the least-squares problem
\[
\arg\min_{q,\{r_{T_i}\}_{i=1}^{12}}
\sum_{i=1}^{12}\sum_{j=1}^{10}
\left(
P(T_i,K_j^i)
-
P^{PC}(C(T_i,K_j^i),q,r_{T_i})
\right)^2,
\]
where $P(T_i,K_j^i)$ denotes the quoted mid-price of a put option with maturity $T_i$ and strike
$K_j^i$, while $C(T_i,K_j^i)$ denotes the corresponding call mid-price. The function $P^{PC}$
gives the put price implied by put-call parity:
\[
P^{PC}(C(T,K),q,r)
=
C(T,K)+Ke^{-rT}-S_0e^{-qT}.
\]
Having inferred the interest rates and dividend yield, we remove in-the-money (ITM) options and
impose the following no-arbitrage restrictions on the remaining observations.
\begin{itemize}
\item We convert put prices into equivalent call prices using put-call parity and, where both call
  and put quotes are available, take the average of the two call prices.
\item We remove observations that violate the no-arbitrage conditions outlined in
  \parencite{Kokholm2016} (Equations 30 to 32). For a fixed maturity, these conditions are
\begin{align*}
&C_t(S_t;T,K_j)
\geq
\max\left\{
0,
S_t e^{-\int_t^T q_s\,\mathrm{d}s}
-
K_j e^{-\int_t^T r_s\,\mathrm{d}s}
\right\},
\tag{Price bounds},\\
&\frac{C_t(S_t;T,K_{j-1})-C_t(S_t;T,K_j)}
{K_j-K_{j-1}}
\in
\left[0,e^{-\int_t^T r_s\,\mathrm{d}s}\right],
\tag{Monotonicity},\\
&\frac{C_t(S_t;T,K_{j-1})-C_t(S_t;T,K_j)}
{K_j-K_{j-1}}
-
\frac{C_t(S_t;T,K_j)-C_t(S_t;T,K_{j+1})}
{K_{j+1}-K_j}
\geq 0,
\tag{Convexity},
\end{align*}
where $K_{j-1}<K_j<K_{j+1}$.

For the monotonicity and convexity conditions, each violation involves two and three observations,
respectively. We iteratively remove the observation involved in the most violations until none
remain. Ties are broken by removing the observation with the widest bid-ask spread, or the first
candidate if bid-ask quotes are unavailable.

\item Finally, we only retain maturities with more than 10 available quotes after filtering.
\end{itemize}

An outline of how much data was removed in each step is included in \Cref{tab:option-filtering}
\begin{table}
  \centering
  \begin{tabular}{lr}
    \toprule
    \textbf{Filtering step} & \textbf{Observations} \\
    \midrule
    Initial observations & 3,778 \\
    Observations after cleaning & 3,665 \\
    \midrule
    Removed ITM options & 1,889 \\
    \quad Original call observations & 417 \\
    \quad Puts converted to calls & 1,359 \\
    Total call observations & 1,776 \\
    \midrule
    Removed by price bounds & 0 \\
    Removed by slope/convexity & 476 \\
    \quad Monotonicity violations found & 11 \\
    \quad Convexity violations found & 7,470 \\
    Removed by maturity filter & 0 \\
    \midrule
    \textbf{Remaining observations} & \textbf{1,300} \\
    Remaining maturities & 12 \\
    \bottomrule
  \end{tabular}
  \caption{Summary of option data filtering.}
  \label{tab:option-filtering}
\end{table}
The inferred dividend yields and interest rates are shown in \Cref{tab:inferred_r_q}
\begin{table}
    \centering
    \begin{tabular}{lr}
        \toprule
        Maturity & Discount Rate (\%) \\
        \midrule
        2019-12-20 & 1.6363 \\
        2020-01-17 & 2.2618 \\
        2020-02-21 & 2.1114 \\
        2020-03-20 & 2.0377 \\
        2020-04-17 & 2.0842 \\
        2020-06-19 & 1.9725 \\
        2020-09-18 & 1.9201 \\
        2020-10-16 & 1.9552 \\
        2020-11-20 & 1.9100 \\
        2020-12-18 & 1.8770 \\
        2021-06-18 & 1.8524 \\
        2021-12-17 & 1.8359 \\
        \midrule
        Inferred dividend yield & 1.8909 \\
        \bottomrule
    \end{tabular}
    \caption{Inferred discount rates by maturity}
    \label{tab:inferred_r_q}
\end{table}
All of this data processing is implemented in \verb|data_handler.py|. 
A Python implementations of the AJD and SVJJ models, together with the Fourier pricing methods
discussed in \Cref{ch:affine,ch:svjj}, is available in \verb|Models/AJDModel.py|,
\verb|Models/SVJJModel.py| and \verb|Fourier.py|, respectively.

The SVJJ model is calibrated jointly to all available observations across the 12 maturities, i.e. we estimate $\theta =
(\lambda,\rho,\overline{v},\kappa_v,\sigma_v, \mu_s,\sigma,\rho_j,\mu_v,V_0),$. 
The initial variance $V_0$ is included as an additional calibration parameter, while the risk-free
interest rate is allowed to vary across maturities according to the previously inferred spot rates.

The calibration is performed by minimizing relative pricing errors. For an option with maturity
$T_i$ and strike $K_j^i$, we define the normalized price error:
\[
\epsilon^i(\theta)
=
\frac{
C_{\mathrm{model}}(T_i,K_j^i;\theta)
-
C_{\mathrm{market}}(T_i,K_j^i)
}{
C_{\mathrm{market}}(T_i,K_j^i)
}.
\]
By minimizing over $\epsilon_j$ we stop options with higher prices from dominating the objective
function. The calibration is implemented in \verb|Calibration/SVJJ_calibration.py|. The initial parameter guesses and
optimization bounds are reported in \Cref{tab:parameter-bounds}, the bounds are chosen to satisfy
the parameter restrictions derived in \Cref{ch:svjj}.
\begin{table}
    \centering
    \begin{tabular}{lrrrrrrrrrr}
        \hline
        & $\lambda$ & $\rho$ & $\overline{v}$ & $\kappa_v$ & $\sigma_v$ & $\mu_s$ & $\sigma$ & $\rho_j$ & $\mu_v$ & $V_0$\\
        \hline
      Initial Guess & 0.100 & -0.200 & 0.040 & 1.000 & 0.300 & -0.100 & 0.150 & -0.100 & 0.200 & 0.020\\
      Lower Bound   & 0.001 & -0.990 & 0.001 & 0.001 & 0.001 & -10.00 & 0.001 & -10.00 & 0.001 &
      0.001 \\
      Upper Bound   & 10.00 & 0.990 & 1.000 & 10.00 & 10.00 & 10.00 & 1.000 & 10.00 & 10.00 &
      10.000\\
        \hline
    \end{tabular}
    \caption{Initial parameter guesses and optimization bounds}
    \label{tab:parameter-bounds}
\end{table}



